\section{Generativity, De-accretion, and Productive Supersession}
\label{sec:exit}

Graduation is not the end of development. A mature productive relation can later disappear because of technological change, higher wages, new competitors, changing comparative advantage, relocation, or deliberate movement into a different productive configuration. The developmental effect of that exit cannot be inferred from the exiting firm alone.

\subsection{Generativity}

\begin{definition}[Generativity]
A graduated relation $e$ is generative for relation $f$ if its operation or accumulated capability improves $f$'s graduation conditions, for example by raising $\Lambda_f$, lowering $q_f^*$, lowering $\delta_f$, or reducing an activation threshold.
\end{definition}

In the positive-complementarity baseline, generativity can expand the graduation frontier:
\begin{equation}
e\text{ graduates}
\rightarrow
\gamma_f\uparrow
\rightarrow
f\in\G
\rightarrow
f\text{ becomes graduable}.
\end{equation}
The propagated object is not merely output. It is improved capacity to create further unsupported productive relations.

\subsection{A simple de-accretion mechanism}

Now consider the reverse direction. Let an upstream supplier $s$ require customer demand above a minimum scale $D_s^*$. Write
\begin{equation}
M_s(x)
=
\bar D_s+\sum_j b_{sj}x_j,
\qquad b_{sj}\ge0.
\label{eq:supplier-demand}
\end{equation}

\begin{proposition}[Direct de-accretion trigger]
\label{prop:deaccretion}
Suppose relation $e$ is active immediately before exit and supplier $s$ satisfies
\[
M_s(x^- )\ge D_s^*,
\]
but
\[
M_s(x^-)-b_{se}<D_s^*.
\]
Then the exit of $e$ makes $s$ nonviable under the minimum-demand condition.
\end{proposition}

\begin{proof}
The exit changes $x_e$ from one to zero and therefore reduces $M_s$ by exactly $b_{se}$. The second inequality places the resulting demand below the viability threshold.
\end{proof}

The proposition is elementary, but its role is important. A firm exit can be a network event rather than an isolated reallocation. If the lost supplier $s$ was itself a complement for other relations, its disappearance can reduce their $\Lambda_f$ or raise their $q_f^*$.

\begin{proposition}[Frontier contraction under complement loss]
\label{prop:frontier-contraction}
Suppose environments $\Omega'$ and $\Omega$ are ordered so that for every relation $f$,
\[
\Lambda_f(\Omega')\le\Lambda_f(\Omega),
\qquad
q_f^*(\Omega')\ge q_f^*(\Omega),
\]
with unchanged $\delta_f$. Then
\[
\G(\Omega')\subseteq\G(\Omega).
\]
\end{proposition}

\begin{proof}
For each $f$, the numerator of the graduation margin weakly falls and the denominator weakly rises, so $\gamma_f(\Omega')\le\gamma_f(\Omega)$. Any relation satisfying $\gamma_f(\Omega')>1$ therefore also satisfies $\gamma_f(\Omega)>1$.
\end{proof}

This yields a useful distinction. An economy can obtain a static allocative gain from replacing an expensive domestic product with a cheaper import while simultaneously losing productive complements that shrink the future graduation frontier. The framework does not determine the welfare sign of that trade-off. It makes the two effects visible so they are not collapsed into one number.

\subsection{Capability after exit}

To represent structural replacement, decompose capability only at the exit stage:
\begin{equation}
q_e=(q_e^S,q_e^R),
\end{equation}
where $q_e^S$ is configuration-specific capability and $q_e^R$ is recombinable capability. The decomposition is intentionally delayed so the baseline graduation theorem remains one-dimensional.

Let $\kappa_e\in[0,1]$ be the fraction of $q_e^R$ released rather than destroyed when $e$ exits, and let $\mu_{ef}\in[0,1]$ be the absorptive transfer coefficient from $e$ into relation $f$. A simple accounting restriction is
\begin{equation}
\sum_f\mu_{ef}\le1.
\end{equation}
The transfer into $f$ is
\begin{equation}
\Delta q_f^R
=
\mu_{ef}\kappa_e q_e^R.
\label{eq:transfer}
\end{equation}

\begin{definition}[Productive supersession]
The exit of $e$ generates productive supersession into $f$ when released recombinable capability from $e$ raises the productive state of $f$ through an identifiable transfer channel.
\end{definition}

\begin{definition}[Successful supersession]
The supersession is successful when the transferred capability contributes to a post-exit productive configuration in which the successor relation is viable without extraordinary support. A sufficient reduced-form condition is
\begin{equation}
q_f^-<q_f^*
\le
q_f^-+\mu_{ef}\kappa_e q_e^R,
\label{eq:supersession-condition}
\end{equation}
combined with satisfaction of the other support-free activation and dependency conditions.
\end{definition}

The empirical burden in this definition is substantial. It is not enough to observe that one industry declined and another later expanded. Actual transfer channels must be identified: workers, engineering practices, supplier competence, machinery, managerial routines, standards, finance, export relationships, or other capabilities. The framework explicitly rejects a story in which Korean shoe capability is simply asserted to have ``become'' semiconductor capability without evidence.

The resulting exit taxonomy is summarized in Table~\ref{tab:exit}.

\begin{table}[ht]
\centering
\caption{Exit modes and capability consequences}
\label{tab:exit}
\small
\begin{tabular}{>{\raggedright\arraybackslash}p{0.22\textwidth}>{\raggedright\arraybackslash}p{0.31\textwidth}>{\raggedright\arraybackslash}p{0.36\textwidth}}
\toprule
\textbf{Exit mode} & \textbf{Immediate mechanism} & \textbf{Developmental diagnostic}\\
\midrule
Failed graduation & Support removed before unsupported viability & Capability stock never crosses the support-free threshold\\
Destructive de-accretion & Exit makes linked relations nonviable & Productive field and possibly the graduation frontier contract\\
Competitive displacement & Rival configuration becomes more attractive & Sign depends on capability destruction versus retention/reallocation\\
Successful supersession & Original activity exits while capability is absorbed elsewhere & Successor relations preserve or deepen unsupported productive capability\\
\bottomrule
\end{tabular}
\end{table}

Two statements follow immediately:
\begin{quote}
\textbf{Industrial death is not capability death.}
\end{quote}
and
\begin{quote}
\textbf{Industrial survival is not developmental success.}
\end{quote}
A permanently protected obsolete activity may survive while trapping resources and capability in a configuration that no longer generates future productive possibilities. Conversely, a declining industry can be part of successful structural transformation if important capabilities persist and recombine.
